% Propietario conservado: /Users/ruben/Documents/ChatGPT/jueces y controles/REVISION_ARGUMENTAL_APERTURAS_CIERRES_20260915/delivery/ARTICLE_V_ES_ARGUMENTAL_20260915/manuscrito/sections/v_entropia_sectorial.tex
% RESULTADO_RECUPERADO desde II REV09. Véase MAPA_PROCEDENCIA.json.
% Selección de líneas y cambios mecánicos explícitos; ninguna prueba nueva.
\subsection{Équation d'état de l'information sectorielle}
\label{v:v:ant:sec:ii-ecuacion-informacion}

L'incidence déjà construite fournit dix-huit facteurs tritiques dans
chaque secteur. Cette multiplicité permet d'écrire explicitement la
relation entre l'aire, la distribution des occupations et l'entropie de
la purification précédente. Nous conservons l'état réduit qui y est défini
et abrégeons $\Delta=\Delta_{\rm mod}$. Pour $m\in\{90,120\}$, nous posons
\begin{equation}
 x_m=e^{-m\Delta},\qquad z_m=1+x_m+x_m^2,\qquad
 f_m=\frac{x_m+2x_m^2}{z_m},\qquad
 v_m=\frac{x_m+4x_m^2}{z_m}-f_m^2.
 \label{v:v:ant:eq:ii-sector-momentos}
\end{equation}
Ici, $f_m$ et $v_m$ sont respectivement la moyenne et la variance de
l'occupation d'un facteur du secteur $m$. Le paramètre $\Delta$ conserve
le statut et la normalisation de l'état réduit préalablement
spécifié.

\subsubsection{Purification quantique bipartite des secteurs d’incidence}
\label{v:v:ant:sec:ii-area-purificacion-producto}

La décomposition de Schmidt précédente se concrétise sur chaque lien.
Pour un poids réel fini \(w\), nous définissons
\[
 z(w)=1+e^{-w}+e^{-2w},\qquad
 \rho(w)=\frac1{z(w)}\sum_{n=0}^2e^{-wn}|n\rangle\langle n|.
\]
Les poids sectoriels sont \(w_m=m\Delta\), avec \(m=90,120\) ;
ainsi, \(z(w_m)=z_m\).

\begin{proposition}[Purification par double thermochamp de l'incidence sectorielle]
\label{v:v:ant:prop:ii-area-purificacion-producto}
Le vecteur normalisé
\begin{equation}
 |\operatorname{TFD}(w)\rangle
 =\frac1{\sqrt{z(w)}}\sum_{n=0}^2
       e^{-wn/2}|n\rangle_A\otimes|n\rangle_{\bar A}
 \label{v:v:ant:eq:ii-area-tfd-enlace}
\end{equation}
purifie \(\rho(w)\). Le produit
\begin{equation}
 |\Psi_\Delta\rangle
 =|\operatorname{TFD}(90\Delta)\rangle^{\otimes18}
  \otimes|\operatorname{TFD}(120\Delta)\rangle^{\otimes18}
 \label{v:v:ant:eq:ii-area-tfd-producto}
\end{equation}
a pour réduction \(\rho_A(\Delta)\) sur les trente-six liens
originaux. Son entropie d'intrication pour la bipartition
\(A\mid\bar A\) est \(s(\Delta)\).
\end{proposition}
\begin{proof}
Le carré de la norme est \(z(w)^{-1}\sum_{n=0}^2e^{-wn}=1\).
Dans la trace partielle, l'orthogonalité du deuxième facteur élimine les
termes d'indices distincts ; les termes restants forment \(\rho(w)\).
La trace partielle commute avec le produit tensoriel. On obtient ainsi
\(\rho(90\Delta)^{\otimes18}\otimes\rho(120\Delta)^{\otimes18}\),
qui coïncide avec l'état normalisé défini précédemment.
L'égalité entre l'entropie d'une réduction et l'intrication de
cette purification se lit directement dans ses coefficients de Schmidt.
\end{proof}

Les deux facteurs de la bipartition sont déterminés au sein de la
construction. La réalisation physique du deuxième facteur conserve cette
identification de sous-système et son algèbre d'observables.

\subsubsection{Moments d'occupation et contenu informationnel}

\begin{proposition}[Aire moyenne et entropie de la purification sectorielle]
Dans cette famille finie d'états,
\begin{align}
 \langle N\rangle_\Delta&=18(f_{90}+f_{120}),\\
 \langle D\rangle_\Delta&=18(90f_{90}+120f_{120}),\\
 \frac{\langle\widehat{\mathcal A}\rangle_\Delta}{\ell_P^2}
 &=18\gamma_{\HMT}a_0(f_{90}+f_{120}),
 \label{v:v:ant:eq:ii-area-media}\\
 s(\Delta)&=18\log z_{90}+18\log z_{120}
                  +\Delta\langle D\rangle_\Delta.
 \label{v:v:ant:eq:ii-entropia-sectorial}
\end{align}
L'entropie dimensionnelle est $S(\Delta)=k_Bs(\Delta)$, avec le
transducteur thermique de la section~\ref{v:v:ant:sec:ii-kb-aut-desarrollo}.
\end{proposition}
\begin{proof}
Dans chaque facteur, les poids normalisés sont
$(1,x_m,x_m^2)/z_m$. Leur premier moment donne $f_m$ et leur deuxième
moment donne la première fraction de $v_m$. La factorisation tensorielle
et l'additivité des opérateurs nombre produisent les trois premières
identités. Enfin,
$\log\rho_A=-(\log Z)I-\Delta D$, avec
$Z=z_{90}^{18}z_{120}^{18}$. Prendre la trace contre $-\rho_A$
donne~\eqref{v:v:ant:eq:ii-entropia-sectorial}. La purification spécifiée
précédemment possède cette même entropie d'intrication.
\end{proof}

La perte d'uniformité possède une lecture quantitative propre.
En dimension \(d_\partial=3^{36}\), soit
\(\tau_\partial=I/d_\partial\). Nous définissons
\begin{equation}
 I_{\rm or}(\Delta)=36\log3-s(\Delta)
 =D\bigl(\rho_A(\Delta)\Vert\tau_\partial\bigr)\ge0.
 \label{v:v:ant:eq:ii-area-deficit-informacional}
\end{equation}
L'égalité s'obtient en substituant
\(\log\tau_\partial=-36\log3\,I\). Si \(p_\nu\) sont les
valeurs propres de \(\rho_A\), la convexité de \(t\log t\) donne
\(\sum_\nu p_\nu\log(d_\partial p_\nu)\ge0\), avec égalité
exactement pour l'état uniforme. Dans cette famille, cela se produit pour
\(\Delta=0\). Le déficit est une information adimensionnelle ; le quantum
aréal \(\gamma_{\HMT}a_0\) est une autre lecture. Chaque expression conserve
sa définition même si leurs évaluations peuvent être proches.

\begin{proposition}[Réponse informationnelle et fluctuations de bord]
Pour tout $\Delta$ réel fini,
\begin{align}
 \frac{\dd\langle D\rangle_\Delta}{\dd\Delta}
 &=-\operatorname{Var}_\Delta(D)
   =-18(90^2v_{90}+120^2v_{120}),\\
 \frac{\dd s}{\dd\Delta}
 &=-\Delta\operatorname{Var}_\Delta(D),
 \label{v:v:ant:eq:ii-respuesta-entropica}\\
 \frac{\dd\langle\widehat{\mathcal A}\rangle_\Delta}{\dd\Delta}
 &=-18\gamma_{\HMT}a_0\ell_P^2(90v_{90}+120v_{120}).
 \label{v:v:ant:eq:ii-respuesta-areal}
\end{align}
\end{proposition}
\begin{proof}
La somme définissant $Z$ est finie. Sa dérivation donne
$\partial_\Delta\log Z=-\langle D\rangle_\Delta$ et
$\partial_\Delta^2\log Z=\operatorname{Var}_\Delta(D)$.
Les secteurs et leurs facteurs sont indépendants dans l'état produit ;
leurs variances s'additionnent donc. En dérivant
\eqref{v:v:ant:eq:ii-entropia-sectorial}, les termes
$-\langle D\rangle_\Delta$ et $\langle D\rangle_\Delta$
s'annulent. Pour l'aire, on utilise
$\partial_\Delta f_m=-mv_m$ dans~\eqref{v:v:ant:eq:ii-area-media}.
\end{proof}

Ces équations relient trois niveaux d’une même construction.
L’incidence fixe les multiplicités et les poids sectoriels ;
l’état réduit fixe leurs probabilités ; les fluctuations de ces
occupations déterminent la réponse de l’aire et de l’entropie.
L’échelle $\gamma_{\HMT}a_0\ell_P^2$ convertit l’occupation en aire,
tandis que $k_B$ convertit l’information en entropie dimensionnelle.
Dans la section~\ref{v:v:ant:sec:gravity}, $\ell_P^2=G\hbar/c^3$ explicite
la dépendance de cette unité d’aire à l’égard de l’action et de la gravitation.
La famille d’états réduits et sa bipartition restent spécifiées
dans ce développement ; toute réalisation d’horizon doit également transporter
ces deux données.

\subsection{Inversion des occupations et invariants entropiques}
\label{v:v:ant:sec:ii-inversion-sectorial}

Soit $\mathcal R$ l'involution unitaire qui échange $|0\rangle$
et $|2\rangle$ et fixe $|1\rangle$ dans chacun des trente-six
facteurs tritiques. Son action sur les occupations conserve toute la
configuration microscopique par un changement bijectif de ses étiquettes.

\begin{theorem}[Entropie conservée et aire complémentaire]
L'involution précédente satisfait
\begin{align}
 \mathcal R N\mathcal R^*&=72I-N,&
 \mathcal R D\mathcal R^*&=7560I-D,\\
 \rho_A(-\Delta)&=\mathcal R\rho_A(\Delta)\mathcal R^*,&
 s(-\Delta)&=s(\Delta),
 \label{v:v:ant:eq:ii-inversion-entropia}\\
 \langle\widehat{\mathcal A}\rangle_{-\Delta}
 +\langle\widehat{\mathcal A}\rangle_\Delta
 &=72\gamma_{\HMT}a_0\ell_P^2.&&
 \label{v:v:ant:eq:ii-area-complementaria}
\end{align}
L'entropie relative entre les deux orientations est
\begin{equation}
 D\bigl(\rho_A(\Delta)\Vert\rho_A(-\Delta)\bigr)
 =\Delta\bigl(7560-2\langle D\rangle_\Delta\bigr).
 \label{v:v:ant:eq:ii-distancia-orientaciones}
\end{equation}
\end{theorem}
\begin{proof}
Chaque occupation locale $n$ se transforme en $2-n$. Ainsi,
$N_m\mapsto36I-N_m$ et
$D\mapsto36(90+120)I-D=7560I-D$.
En outre, $z_m(-\Delta)=e^{2m\Delta}z_m(\Delta)$,
d'où $Z(-\Delta)=e^{7560\Delta}Z(\Delta)$.
Substituer les deux identités dans l'état normalisé prouve la
conjugaison unitaire. L'invariance spectrale prouve l'égalité
des entropies ; la transformation de $N$ prouve la somme des aires.
Enfin, en soustrayant les deux logarithmes des états et en prenant
leur espérance dans $\rho_A(\Delta)$, on obtient
\eqref{v:v:ant:eq:ii-distancia-orientaciones}.
\end{proof}

Pour $\Delta=0$, l'état réduit est uniforme et
$s(0)=36\log3$. À la limite $\Delta\to+\infty$, survit
la configuration d'occupation nulle ; l'entropie tend vers zéro.
L'inversion mène cette limite à la configuration d'occupation maximale.
Aux deux extrêmes, l'entropie est la même et l'aire distingue les
configurations. Pour $\Delta>0$, la réponse
\eqref{v:v:ant:eq:ii-respuesta-entropica} est strictement négative puisque
l'état est de rang plein et $D$ possède plus d'une valeur propre.
La loi conserve ainsi la différence entre réversibilité de la
configuration complète, égalité des entropies et changement de lecture aréale.
Ces propriétés appartiennent à la famille sectorielle spécifiée ;
le transport physique de la bipartition est traité dans
la section~\ref{v:v:ant:sec:biparticion-transporte}.

\subsection{Première loi modulaire et variations finies d'aire}
\label{v:v:ant:sec:ii-area-ley-modular}

Les dérivées de \eqref{v:v:ant:eq:ii-respuesta-entropica} font varier le paramètre
\(\Delta\) d'une famille d'états. La première loi modulaire fixe,
en revanche, un état de référence et considère les variations de l'état
observé autour de celui-ci. Prenons
\(\rho_0=\rho_A(\Delta_0)\), avec \(\Delta_0\) réel et fini,
et maintenons fixes \(\Delta_0,\gamma_{\HMT},a_0,\ell_P\).
On conserve la branche \(\gamma_{\HMT}>0\), \(a_0>0\).
Les aires sectorielles normalisées sont
\begin{equation}
 \mathcal A_m^{\rm n}=\gamma_{\HMT}a_0N_m,
 \qquad
 \widehat{\mathcal A}_m=\ell_P^2\mathcal A_m^{\rm n},
 \qquad m\in\{90,120\}.
 \label{v:v:ant:eq:ii-area-sector-normalizado}
\end{equation}
Le logarithme de l'état de référence prend la forme
\begin{equation}
 K_0=-\log\rho_0
 =c_0I+\beta_{90}\mathcal A_{90}^{\rm n}
         +\beta_{120}\mathcal A_{120}^{\rm n},
 \qquad
 \beta_m=\frac{m\Delta_0}{\gamma_{\HMT}a_0},
 \quad c_0=\log Z(\Delta_0).
 \label{v:v:ant:eq:ii-area-beta-sectorial}
\end{equation}
En particulier, \(\beta_{120}=\tfrac43\beta_{90}\) ; pour
\(\Delta_0\ne0\), leur rapport est \(4/3\). Les coefficients
résultent de l'expression du même opérateur modulaire en unités aréales.

\begin{theorem}[Première loi modulaire en coordonnées aréales]
\label{v:v:ant:thm:ii-area-primera-ley}
Soit \(\sigma(\varepsilon)\) une famille différentiable d'états
normalisés sur un intervalle contenant zéro, avec
\(\sigma(0)=\rho_0\). Alors
\begin{equation}
 \left.\frac{\dd s(\sigma(\varepsilon))}{\dd\varepsilon}\right|_0
 =\sum_{m=90,120}\beta_m
   \left.\frac{\dd\langle\mathcal A_m^{\rm n}\rangle_
                       {\sigma(\varepsilon)}}{\dd\varepsilon}\right|_0.
 \label{v:v:ant:eq:ii-area-primera-ley}
\end{equation}
\end{theorem}
\begin{proof}
L'état \(\rho_0\) est de rang plein. Dans son voisinage, la dérivée
de la trace de \(-\sigma\log\sigma\) est
\(-\Tr[\dot\sigma(\log\rho_0+I)]\) ; elle s'obtient en diagonalisant
\(\rho_0\) dans la formule de la dérivée spectrale de la trace.
La normalisation implique \(\Tr\dot\sigma=0\), de sorte que la
dérivée est \(\Tr(\dot\sigma K_0)\). Substituer
\eqref{v:v:ant:eq:ii-area-beta-sectorial} annule également \(c_0I\) et
produit l'identité.
\end{proof}

\begin{theorem}[Identité finie d'aire et déficit d'entropie relative]
\label{v:v:ant:thm:ii-area-ley-finita}
Pour tout état normalisé \(\sigma\) du même espace fini,
définissons \(\Delta_\sigma\langle B\rangle
=\Tr[(\sigma-\rho_0)B]\). On a
\begin{equation}
 \boxed{
 s(\sigma)-s(\rho_0)
 =\sum_{m=90,120}\beta_m
      \Delta_\sigma\langle\mathcal A_m^{\rm n}\rangle
   -D(\sigma\Vert\rho_0).}
 \label{v:v:ant:eq:ii-area-ley-finita}
\end{equation}
Ici, \(D(\sigma\Vert\rho_0)=\Tr[\sigma(\log\sigma-\log\rho_0)]\)
est finie, non négative, et s'interprète avec \(0\log0=0\).
\end{theorem}
\begin{proof}
Par définition,
\(D(\sigma\Vert\rho_0)=-s(\sigma)+\Tr(\sigma K_0)\), tandis que
\(s(\rho_0)=\Tr(\rho_0K_0)\). La soustraction donne
\(s(\sigma)-s(\rho_0)=\Delta_\sigma\langle K_0\rangle
-D(\sigma\Vert\rho_0)\). Substituer
\eqref{v:v:ant:eq:ii-area-beta-sectorial} démontre la formule puisque les
traces des deux états valent un.

Pour voir la positivité, diagonalisons \(\sigma\) avec les valeurs propres
\(p_i\) et les vecteurs \(u_i\), et \(\rho_0\) avec les valeurs propres
\(q_j>0\) et les vecteurs \(v_j\). Posons
\(r_i=\langle u_i,\rho_0u_i\rangle\). La concavité du logarithme
fournit
\(\sum_j|\langle u_i,v_j\rangle|^2\log q_j\le\log r_i\).
D'où
\(D(\sigma\Vert\rho_0)\ge\sum_i p_i\log(p_i/r_i)\ge0\) ;
la dernière inégalité est la convexité de \(t\log t\) appliquée
avec les poids \(r_i\), de somme un. La positivité stricte des
\(q_j\) garantit que tous les logarithmes de référence sont finis.
\end{proof}

L'identité finie conserve la distance informationnelle par rapport à
l'état de référence ; son terme tangent est
\eqref{v:v:ant:eq:ii-area-primera-ley}. Pour les aires physiques, les
coefficients correspondants sont \(\beta_m/\ell_P^2\).
Multiplier l'équation par le transducteur commun \(k_B\) publie
l'entropie dimensionnelle. Ainsi, la relation entre géométrie et information
demeure exacte également pour les variations finies, avec sa correction
explicite et sans modifier les paramètres de l'état de référence.
